\documentclass[11pt,a4paper]{article}
\usepackage{curso}
\begin{document}
\portada{Lección}{0.6}{Análisis de Fourier en grupos abelianos finitos}

\noindent\textbf{Unidad 0.}\quad\textbf{Objetivos:} definir caracteres de $\Z/r\Z$; probar su ortogonalidad;
usar la DFT para diagonalizar caminatas en el ciclo. \textbf{Prerrequisitos:} números complejos; lección 0.5.

\section{¿Por qué importa?}
La Unidad 7 compara el algoritmo de Shor con un algoritmo termodinámico de búsqueda de orden por difusión
\cite{cadavid1,cadavid2}. Ambos son, en el fondo, análisis de Fourier sobre $\Z/r\Z$: Shor amplifica picos de
Fourier por interferencia; la difusión amortigua todos los modos excepto el trivial y lee el período en la
\emph{probabilidad de regreso}.

\section{Caracteres}
\begin{definicion}[ (carácter)]
Un carácter de un grupo abeliano finito $G$ es un homomorfismo $\chi:G\to\{z\in\mathbb C:|z|=1\}$,
$\chi(x+y)=\chi(x)\chi(y)$. Para $G=\Z/r\Z$: $\chi_k(x)=e^{2\pi ikx/r}$, $k=0,\dots,r-1$.
\end{definicion}
Interpretación: $\chi_k$ envía la hora $x$ del reloj de $r$ horas al punto del círculo unitario a $kx/r$ vueltas;
sumar horas es componer rotaciones. Hay exactamente $r$ caracteres (el grupo dual $\hat G\cong G$), y
$\chi_k\chi_l=\chi_{k+l}$, $\overline{\chi_k}=\chi_{-k}$.

\section{Ortogonalidad}
\begin{lema}
$\displaystyle\sum_{x=0}^{r-1}e^{2\pi ixk/r}=\begin{cases}r,&r\mid k\\0,&\text{si no.}\end{cases}$
\end{lema}
\begin{proof}
Si $r\mid k$ cada término vale 1. Si no, $\omega=e^{2\pi ik/r}\neq1$ y la suma geométrica es
$\frac{\omega^r-1}{\omega-1}=0$ porque $\omega^r=1$. Geométricamente: los vectores puestos punta con cola forman un
polígono regular (recorrido $\gcd(k,r)$ veces) que se cierra.
\end{proof}
\begin{resultado}[ (ortogonalidad)]
$\langle\chi_k,\chi_l\rangle=\sum_x\chi_k(x)\overline{\chi_l(x)}=\sum_x\chi_{k-l}(x)=r\,\delta_{kl}$.
\end{resultado}

\section{Transformada de Fourier discreta}
Los $r$ caracteres son ortogonales en $\mathbb C^r$, luego forman una base:
\begin{equation}\hat f(k)=\sum_xf(x)\overline{\chi_k(x)},\qquad f(x)=\frac1r\sum_k\hat f(k)\chi_k(x),\qquad
\sum_x|f(x)|^2=\frac1r\sum_k|\hat f(k)|^2.\end{equation}
\begin{ejemplo}
$r=4$, $f=(1,2,0,1)$, $\chi_1(x)=i^x$: $\hat f(0)=4$, $\hat f(1)=1-2i+0+i=1-i$, $\hat f(2)=1-2+0-1=-2$,
$\hat f(3)=1+i$. Parseval: $\sum|f|^2=6=\frac14(16+2+4+2)$.
\end{ejemplo}
Convolución: $(f*g)(x)=\sum_yf(y)g(x-y)$ cumple $\widehat{f*g}=\hat f\,\hat g$.

\section{Caminatas aleatorias y núcleo de calor}
Una caminata en $\Z/r\Z$ con distribución de paso $\mu$ (probabilidad de moverse $s$) actúa sobre distribuciones
por convolución: $p_{n+1}=p_n*\mu$. Por tanto $\hat p_n(k)=\hat p_0(k)\,\hat\mu(k)^n$ y los caracteres son vectores
propios con valores propios $\lambda_k=\overline{\hat\mu(k)}$ (reales si $\mu$ es simétrica). Para la caminata
perezosa ($\mu(0)=\frac12$, $\mu(\pm1)=\frac14$): $\lambda_k=\frac12\big(1+\cos\frac{2\pi k}r\big)\in[0,1]$. Partiendo de $0$:
\begin{equation}p_n(x)=\frac1r\sum_k\lambda_k^n\chi_k(x),\qquad p_n(0)=\frac1r\sum_{k=0}^{r-1}\lambda_k^n .\end{equation}
Este es el \emph{núcleo de calor discreto}. Cada modo decae como $\lambda_k^n$; solo $k=0$ ($\lambda_0=1$) sobrevive y
$p_n\to$ uniforme. La brecha espectral es $1-\lambda_1=\frac12(1-\cos\frac{2\pi}r)\approx\pi^2/r^2$.

\begin{aplicacion}
La búsqueda de orden por difusión (Unidad 7.4) usa una caminata con generadores $b^{\pm2^t}$ y valores propios
$\lambda_k=\frac12\big(1+\frac1{M+1}\sum_t\cos\frac{2\pi k2^t}{r}\big)$; la probabilidad de regreso
$p_n(e)=\frac1r\sum_k\lambda_k^n\to\frac1r$, de modo que $r=\mathrm{round}(1/p_n(e))$ para $n$ suficientemente grande.
Las identidades de sumas trigonométricas de la lección 7.5 provienen de calcular $p_n$ de dos maneras \cite{cadavid3}.
\end{aplicacion}

\section{Grupos abelianos finitos en general}
Todo grupo abeliano finito es isomorfo a $\Z/r_1\times\dots\times\Z/r_m$; sus caracteres son
$\chi_{\mathbf k}(\mathbf x)=\prod_je^{2\pi ik_jx_j/r_j}$ y toda la teoría anterior se extiende componente a componente.

\begin{error}
\begin{itemize}[nosep]
\item Conjugado en el producto interno; factor $1/r$ en la inversa (o $1/\sqrt r$ en ambas, convención unitaria).
\item Frecuencia $k$ y posición $x$ viven en copias distintas de $\Z/r\Z$.
\item La suma de raíces vale $r$ cuando $r\mid k$, no solo para $k=0$.
\end{itemize}
\end{error}

\section{Resumen}
$\chi_k(x)=e^{2\pi ikx/r}$; ortogonalidad $r\delta_{kl}$; DFT e inversa; las caminatas se diagonalizan y
$p_n(0)=\frac1r\sum_k\lambda_k^n$.

\begin{thebibliography}{9}\small
\bibitem[32]{cadavid1} C. A. Cadavid et al., An integer factorization algorithm which uses diffusion as a computational engine, 2023.
\bibitem[33]{cadavid2} C. A. Cadavid et al., Diffusion computation versus quantum computation, 2026.
\bibitem[34]{cadavid3} C. A. Cadavid et al., On an approach for evaluating certain trigonometric character sums using the discrete time heat kernel, 2022.
\bibitem[4]{nc} M. A. Nielsen e I. L. Chuang, \emph{Quantum Computation and Quantum Information}, 2010.
\end{thebibliography}
\end{document}
